% !TeX root = ../main.tex
\section{灾害前环境异常与关键关联因素识别（问题一）}\label{sec:factors}

\subsection{基于灾害机制的候选特征构建}\label{sec:candidate-features}
% 依据物理机制和文献，将原始观测转化为正积温、累计降水、净辐射、土壤湿润指数、湍流热交换等候选特征。
\writingplaceholder

\subsection{历史背景与气象异常识别}\label{sec:anomaly-detection}
% 构建季节和昼夜变化基线，识别灾前不同时间窗口的异常。
\writingplaceholder

\subsection{环境特征与灾害事件的关联分析}\label{sec:association}
% 利用事件与历史对照、滞后分析等方法研究关联性。若能够补充多次灾害标签，可进一步采用统计或机器学习方法。
\writingplaceholder

\subsection{关键因素排序与稳健性检验}\label{sec:factor-ranking}
% 输出题目要求的前五个候选关联因素，分析排名对特征定义、对照时段和时间窗口的敏感性。
\writingplaceholder

\subsection{气象异常的物理解释}\label{sec:physical-interpretation}
% 讨论温度、降水、辐射等异常与冰岩崩、泥石流灾害链之间可能存在的联系及证据局限。
\writingplaceholder
